(a) Fix $P\in Y$ and choose another point $Q\in Y$. Since $d>1$, the line $PQ$ is not $Y$. By \exref{I.5.4},
\[d=(PQ\cdot Y)\geq(PQ\cdot Y)_P+(PQ\cdot Y)_Q\geq\mu_P(Y)+1.\]
Thus $\mu_P(Y)\leq d-1$.

(b) Put the point of multiplicity $d-1$ at $(0:0:1)$. The affine equation is $f_{d-1}(x,y)+f_d(x,y)=0$, with both homogeneous terms nonzero. They are relatively prime, since a common factor would factor the equation of $Y$. Projection from the origin is $(x,y)\mapsto t=x/y$. On the open set where the following expressions are defined and nonzero, its inverse is
\[t\mapsto\left(-\frac{t f_{d-1}(t,1)}{f_d(t,1)},-\frac{f_{d-1}(t,1)}{f_d(t,1)}\right).\]
Indeed, division of the equation by $y^{d-1}$ gives $f_{d-1}(t,1)+y f_d(t,1)=0$. Thus $Y$ is birational to $\mathbb{P}^1$.
